Facial-expression recognition (FER) often assigns one category to an image, although annotators may disagree. FER+ retains this variation through multiple responses for FER2013 images \cite{goodfellow2013fer,barsoum2016ferplus}. Training on those distributions can improve agreement with human response frequencies without improving categorical decisions \cite{washington2021crowdsourcing}. Likewise, soft-label training does not necessarily improve selective prediction, in which a model withholds its less confident decisions \cite{khurana2024crowd,xia2025smoothing}. These distinctions are established; our question is what a controlled FER comparison adds to their interpretation.

Three choices complicate that comparison. First, replacing a hard target with a vote distribution changes its largest probability, treatment of tied maxima, and minority-category allocation together. A hard--soft difference cannot attribute the result to one of these changes. Second, models at equal coverage---the accepted fraction of images---may accept different images. Their reported disagreement combines category choices with the composition of those sets. Third, annotations can include categories absent from the model. Removing those votes changes the scoring reference and hides part of the disagreement.

We examine these choices using matched training targets, cross-evaluation on identical accepted images, and complete-vocabulary scoring. The original primary follow-up compares seven-output empirical vote targets (Soft) with a control retaining the largest target probability but spreading the remainder uniformly (Uniform). Tie-preserving controls address alternative handling of tied maxima. Initialization, images, weights, order, augmentation and optimization remain matched within each backbone and seed. The controls preserve specified properties rather than all aspects of the targets: Uniform also changes entropy and aggregate class proportions.

The study addresses three questions:
\begin{enumerate}
\item \textbf{Target information:} What selective advantage remains after controlling dominant target probabilities and ties, and after preserving each target's probability values while disrupting minority-category identities?
\item \textbf{Evaluation choice:} How does predictor ordering depend on accepted images and on whether the model includes seven or eight expression outputs?
\item \textbf{Robustness:} Which findings persist across tested backbones, confidence rankings, checkpoints, calibration transfer and external clips?
\end{enumerate}

Our contribution is an empirical comparison of these dependencies, rather than a new architecture or rejection algorithm. The original seven-output Swin-T comparison uses three paired seeds at 80\% coverage. Secondary analyses examine ResNet-18, fixed-epoch checkpoints and alternative confidence scores; follow-up controls test minority-category allocation and an eight-expression vocabulary. Together, the phases contain 36 training runs and 81 named checkpoint-rule states, not necessarily distinct weight sets. They are adaptive and reuse one image source, with each new comparison specified before its own outcomes. Frozen evaluation on 1,440 RAVDESS clips \cite{livingstone2018ravdess} assesses a different population and prediction unit, providing a transfer test rather than independent static-image confirmation.

A local messaging preview motivated the initial seven-category model. Related systems use text-based emotion notifications, facial and textual signals, or expression-informed emoji suggestions \cite{wang2016emotionpush,chong2019emochat,liu2019selfawareness}. Such cues do not establish a person's internal emotion or intention \cite{barrett2019emotions}. We retain the implementation as a supplementary artifact; the research evidence concerns annotation-based recognition and selection.
